\section{Bilayer spectral decomposition of the angular operator}

Let \(R=R^*=R^{-1}\) be the active involution and

\[
P_\pm=\frac{I\pm R}{2},
\qquad
T=q_+P_+ +q_-P_-,
\qquad 0<q_+<q_-<1.
\]

With
\[
x=A_{\rm rad}=\frac{\pi A}{180},\qquad
y=C^*_{\rm rad}=\frac{\pi C^*}{180},
\]
the two channels are constructed before evaluating the vacuum:
\begin{equation}
q_+=e^{-x-y},\qquad q_-=e^{-x+y}.
\label{eq:vacuum-angular-channels}
\end{equation}
The exchange of sheets reverses \(C^*\), that is, \(y\mapsto-y\), and permutes \(q_+\leftrightarrow q_-\). The pair is thus the spectral reading of the same angular antecedent, not two parameters fitted separately.

\begin{definition}[Constitutive response]
\begin{equation}
\mathcal V_{90,120}=(I-T^{90})(I-T^{120})^{-1}.
\label{eq:vacuum-operator}
\end{equation}
\end{definition}

Since \(\|T\|<1\), the denominator is positive and invertible. Functional calculus gives

\begin{equation}
\mathcal V_{90,120}=r_+P_+ +r_-P_-,
\qquad
r_\pm=\frac{1-q_\pm^{90}}{1-q_\pm^{120}}.
\label{eq:vacuum-eigenvalues}
\end{equation}

The assertion of invertibility deserves to be made explicit. For each leaf,
\(0<q_\pm^{120}<1\); therefore
\[
I-T^{120}=(1-q_+^{120})P_++(1-q_-^{120})P_-
\]
is strictly positive and
\[
(I-T^{120})^{-1}
=\frac{P_+}{1-q_+^{120}}+
 \frac{P_-}{1-q_-^{120}}.
\]
The response \(\mathcal V_{90,120}\) is not obtained by solving four
separate scalar equations: it is the result of the functional calculus of a
single positive contraction. This observation prevents permeability,
permittivity, impedance and propagation from losing their shared antecedent.

\begin{theorem}[Joint spectral determination of the four constitutive relations]
The four constituent characters
\begin{equation}
\widehat\varepsilon=r_+^2,\qquad
\widehat\mu=r_-^2,\qquad
\widehat Z=\frac{r_-}{r_+},\qquad
\widehat c=\frac1{r_+r_-}
\label{eq:vacuum-four}
\end{equation}
are spectral functions of \(\mathcal V_{90,120}\). They satisfy
\begin{equation}
\widehat\mu\widehat\varepsilon=\widehat c^{-2},
\qquad
\frac{\widehat\mu}{\widehat\varepsilon}=\widehat Z^2.
\label{eq:vacuum-relations}
\end{equation}
\end{theorem}

\begin{proof}
The identities are obtained by substituting \cref{eq:vacuum-four}. No dimensional chart nor any exterior value intervenes.
\end{proof}

\begin{corollary}[Reconstruction of the two sheets]
The four characters do not contain four degrees of freedom. Any of
the pairs \((\widehat\mu,\widehat\varepsilon)\),
\((\widehat Z,\widehat c)\), or \((\mathcal E,\mathcal B)\) reconstructs the
two positive eigenvalues:
\begin{align}
r_-&=\sqrt{\widehat\mu}
=\sqrt{\widehat Z/\widehat c},
&r_+&=\sqrt{\widehat\varepsilon}
=\frac1{\sqrt{\widehat Z\widehat c}},
\label{eq:vacuum-reconstruct-r}\\
\log r_-&=\frac{\mathcal E+\mathcal B}{2},
&\log r_+&=\frac{\mathcal E-\mathcal B}{2}.
\label{eq:vacuum-reconstruct-log}
\end{align}
In particular, orientation is lost if only the product is retained,
but is recovered by adding the quotient.
\end{corollary}

\begin{proof}
The first line is the positive root of the identities
\cref{eq:vacuum-four,eq:vacuum-relations}. The second inverts the linear
system
\(\mathcal E=\log r_-+\log r_+\),
\(\mathcal B=\log r_--\log r_+\).
\end{proof}

This corollary expresses a central property of HMT: product and quotient are
complementary functionals of the same state. The product determines the
propagation in common, whereas the quotient preserves the relative orientation
of the inner and outer sheets.

In the V2 certificate,

\begin{align*}
r_+&=0.9999996285482493631786952281\ldots,\\
r_-&=0.9997241371895183641315165853\ldots,\\
\widehat\varepsilon&=0.9999992570966367027604416155\ldots,\\
\widehat\mu&=0.9994483504793269350899815559\ldots,\\
\widehat Z&=0.9997245085389372154555543466\ldots,\\
\widehat c&=1.0002763104861575261063862689\ldots.
\end{align*}

These numerals are the terminal recognition of \cref{eq:vacuum-operator}; they do not define its exponents \(90,120\).

\section{Harmonic factorization of the \(90\) and \(120\) sectors}

The harmonic semigroup is

\[
\langle90,120\rangle=30\langle3,4\rangle.
\]

If \(s=q^{30}\), then

\begin{equation}
\frac{1-q^{90}}{1-q^{120}}
=\frac{1-s^3}{1-s^4}
=\frac{1+s+s^2}{1+s+s^2+s^3}.
\label{eq:three-four-completion}
\end{equation}

The identity also separates the response from its defect. If
\[
F_{3\to4}(s)=\frac{1+s+s^2}{1+s+s^2+s^3},
\qquad 0<s<1,
\]
then
\begin{equation}
d_{3\to4}(s):=1-F_{3\to4}(s)
=\frac{s^3}{1+s+s^2+s^3}>0.
\label{eq:vacuum-defect34}
\end{equation}
The additional term is preserved exactly as a positive
defect. Moreover,
\begin{equation}
F_{3\to4}'(s)
=-\frac{s^2(3+2s+s^2)}{(1+s+s^2+s^3)^2}<0.
\label{eq:vacuum-monotonic34}
\end{equation}
Consequently, the order of the angular channels uniquely determines
the order of the constitutive responses. The assignment of
\(\widehat\mu\) and \(\widehat\varepsilon\) to the respective sheets is obtained
analytically, without resorting to their decimal expansions.

Each sheet of the vacuum constitutes the exact completion from three to four
terms of the same elementary mode \(30\). Consequently, permeability
and permittivity are interpreted as two quadratic functionals of the same
completion evaluated on conjugate sheets, and not as independent
constitutive parameters.

Since the function \((1-q^{90})/(1-q^{120})\) decreases in \(0<q<1\) and \(q_+<q_-\), we obtain

\[
0<r_-<r_+<1,\qquad
0<\widehat\mu<\widehat\varepsilon<1,\qquad
0<\widehat Z<1<\widehat c.
\]

The limits of the functional determine its geometry. When \(q\to0^+\),
\(F(q)\to1\): the four readings approach unity and the memory
defect disappears. When \(q\to1^-\),
\[
F(q)\longrightarrow\frac{90}{120}=\frac34.
\]
The response therefore takes values in the positive interval \((3/4,1)\). The
endpoint \(3/4\) is not a fitted constant: it is the quotient of the two
channel lengths and the continuous realization of the
three--to--four completion.

\section{Sheet involution and preservation of orientation}

The sheet exchange \(C^*\mapsto-C^*\) acts by

\begin{equation}
\widehat\mu\longleftrightarrow\widehat\varepsilon,\qquad
\widehat Z\longmapsto\widehat Z^{-1},
\qquad
\widehat c\longmapsto\widehat c.
\label{eq:vacuum-sheet}
\end{equation}

The common propagation is even; impedance preserves orientation.
This distinction will reappear in Barbero and CKM: not all functionals
associated with the change of sheet have the same parity.

\section{Logarithmic characters and tracial compatibility}

Define

\[
\mathcal E=\log(r_+r_-),\qquad
\mathcal B=\log(r_-/r_+).
\]

With the normalized trace \(\tau=\Tr/2\),

\begin{equation}
2\tau(\log\mathcal V)=\mathcal E,
\qquad
-2\tau(R\log\mathcal V)=\mathcal B.
\label{eq:vacuum-traces}
\end{equation}

The pair \((\mathcal E,\mathcal B)\) constitutes a linear coordinate system for the operator
logarithm:
\begin{equation}
\log\mathcal V
=\frac{\mathcal E}{2}I-\frac{\mathcal B}{2}R.
\label{eq:vacuum-log-decomposition}
\end{equation}
The scalar component is invariant under sheet exchange; the
oriented component changes sign. This decomposition directly
demonstrates the parity of propagation and the oriented character of
impedance.

The nonadic amplification

\[
\mathcal V_m=\mathcal V\otimes I_{9^m}
\]

is compatible with the expectations \(E_m=\mathrm{id}\otimes\tau_9\):

\[
E_m(\mathcal V_{m+1})=\mathcal V_m.
\]

Therefore, \(\mathcal E\) and \(\mathcal B\) are not peculiarities of a \(2\times2\) matrix; they persist in the UHF tower and in its tracial closure \(II_1\).

More precisely, for every polynomial \(p\) and afterwards, by continuity,
for every continuous function on the spectrum,
\begin{equation}
E_m\bigl(p(\mathcal V_{m+1})\bigr)=p(\mathcal V_m).
\label{eq:vacuum-functional-refinement}
\end{equation}
Compatibility does not preserve only two numbers: it preserves the entire
functional algebra generated by the response. In particular, it preserves its
spectral projectors, powers, logarithm, and quadratic forms
that measure the defects \cref{eq:vacuum-defect34}.

\section{Dimensional realization of the constitutive relations}

Dimensional Mechanics determines

\begin{align}
Z_{\rm vac}&=\widehat Z\,u_Su_Q^{-2},\\
c_{\rm vac}&=\widehat c\,u_C,\\
\mu_{\rm vac}&=\widehat\mu\,u_Su_C^{-1}u_Q^{-2},\\
\varepsilon_{\rm vac}&=\widehat\varepsilon\,u_S^{-1}u_C^{-1}u_Q^2.
\label{eq:vacuum-sections}
\end{align}

The constitutive identities hold on the straight lines:

\begin{equation}
\mu_{\rm vac}\varepsilon_{\rm vac}=c_{\rm vac}^{-2},
\qquad
\mu_{\rm vac}/\varepsilon_{\rm vac}=Z_{\rm vac}^2.
\end{equation}

The SI chart comes later. It is not necessary to declare \(c\) primitive and then solve \(\mu\varepsilon=c^{-2}\): the three objects descend jointly from \(\mathcal V\) and from the declared lines.

The separation of levels is expressed by the diagram
\[
\begin{array}{c}
(A,C^*)\longrightarrow T\longrightarrow\mathcal V
\longrightarrow
(\widehat\mu,\widehat\varepsilon,\widehat Z,\widehat c)
\\[2mm]
\Big\downarrow\text{ line charts}\hspace{36mm}
\Big\downarrow\text{ final evaluation}
\\[2mm]
(u_S,u_Q,u_C)\longrightarrow
(\mu_{\rm vac},\varepsilon_{\rm vac},Z_{\rm vac},c_{\rm vac})
\longrightarrow\text{ SI coordinates}.
\end{array}
\]
The first line is dimensionless and operatorial; the second determines types of
quantity; the third selects a metrological chart. Altering the final chart
changes neither the eigenvalues nor the constitutive identities. In this
precise sense, mathematical--physical unification consists in distinguishing
the genealogical invariant from its coordinate system, not in removing the
dimensional structure.

\begin{proposition}[Spectral minimality]
Let \(W\) be a positive two-sheet operator with simple spectrum
\(\{r_+,r_-\}\). If four positive characters satisfy
\(\mu\varepsilon=c^{-2}\) and \(\mu/\varepsilon=Z^2\), if their propagation in common is determined by \(c^{-1}=\det W=r_-r_+\), and if its orientation
fixes \(Z=r_-/r_+\), then necessarily
\[
(\mu,\varepsilon,Z,c)
=\left(r_-^2,r_+^2,\frac{r_-}{r_+},\frac1{r_-r_+}\right).
\]
Therefore, \cref{eq:vacuum-four} is not one parametrization among many:
it is the minimal positive realization compatible with product, quotient and
orientation.
\end{proposition}

\begin{proof}
From \(Z=r_-/r_+\) and \(c^{-1}=\det W=r_-r_+\) the positive roots are obtained uniquely
roots. The two constitutive identities then force afterwards
\(\mu=r_-^2\) and \(\varepsilon=r_+^2\).
\end{proof}

